\clearpage\thispagestyle{empty}
\pdfbookmark[0]{Sources for the review of the Milan ephemerides for 1816}{carlini1816-sources}
\begin{center}\Large\bfseries Sources for the review of the Milan ephemerides for 1816\end{center}

The translation follows Carl Friedrich Gauss’s review in \emph{Werke}, vol.~VI (Göttingen, 1874), pp.~592--593. The collected edition has the page reference \emph{457..46c}, with an open, c-shaped final glyph. The first printing appeared in \emph{Göttingische gelehrte Anzeigen}, no.~47, 23 March 1816, pp.~457--460.

In the passage on the spirit level, the first printing (p.~459) gives $1\frac{2}{3}$ lines. The collected edition has an indistinct denominator. The translation retains its outline. The decimal point in Vesta’s eccentricity and the unclosed parenthesis in the orbital inclination entry also follow the collected edition. The historical units are retained. Cesaris gives the rainfall in Paris inches and lines (appendix p.~88). The forms \emph{nordliche}, \emph{Erfahrnngen}, \emph{Mitagsfernrohr} and \emph{häufigern} in the collected text are rendered according to their meanings.

The title page of Francesco Carlini’s \emph{Effemeridi astronomiche di Milano per l’anno bisestile 1816} gives Milan, Cesarea Regia Stamperia, 1815. The copy digitized by Google and available through the INAF Brera Observatory bears a stamp of the Royal University of Rome. Its contents list the three supplements to the calendar, followed by the appendix. Oriani’s first paper describes a repeating circle three feet in diameter. His summary of the solstice observations is on p.~85 of the appendix.

AI-translated from German. AI-produced typesetting.

\href{https://gdz.sub.uni-goettingen.de/id/PPN236007467}{Collected edition}. First printing, pp.~\href{https://images.p.sub.uni-goettingen.de/iiif/image/gdz:PPN319721507_1816:00000459/full/1650,/0/default.jpg}{457}, \href{https://images.p.sub.uni-goettingen.de/iiif/image/gdz:PPN319721507_1816:00000460/full/1650,/0/default.jpg}{458}, \href{https://images.p.sub.uni-goettingen.de/iiif/image/gdz:PPN319721507_1816:00000461/full/1650,/0/default.jpg}{459}, \href{https://images.p.sub.uni-goettingen.de/iiif/image/gdz:PPN319721507_1816:00000462/full/1650,/0/default.jpg}{460}. \href{https://www1.brera.inaf.it/~mario.carpino/Biblioteca_Digitale/Effemeridi_Milano/mirror/EM1816b.pdf#page=2}{Yearbook title page}. \href{https://www1.brera.inaf.it/~mario.carpino/Biblioteca_Digitale/Effemeridi_Milano/mirror/EM1816b.pdf#page=8}{Contents}. \href{https://www1.brera.inaf.it/~mario.carpino/Biblioteca_Digitale/Effemeridi_Milano/mirror/EM1816b.pdf#page=120}{Oriani on refraction}. \href{https://www1.brera.inaf.it/~mario.carpino/Biblioteca_Digitale/Effemeridi_Milano/mirror/EM1816b.pdf#page=130}{Cesaris on the movement of buildings}. \href{https://www1.brera.inaf.it/~mario.carpino/Biblioteca_Digitale/Effemeridi_Milano/mirror/EM1816b.pdf#page=150}{Carlini’s Vesta table}. \href{https://www1.brera.inaf.it/~mario.carpino/Biblioteca_Digitale/Effemeridi_Milano/mirror/EM1816b.pdf#page=202}{Oriani’s solstice results}. \href{https://www1.brera.inaf.it/~mario.carpino/Biblioteca_Digitale/Effemeridi_Milano/mirror/EM1816b.pdf#page=205}{Cesaris on rainfall}.

\clearpage
